\documentclass{article}
\usepackage[utf8]{inputenc}
\usepackage[T1]{fontenc}
\usepackage{geometry}
\usepackage{graphicx}
\usepackage{float}
\usepackage{hyperref}
\usepackage[english,indonesian]{babel}

\title{Optimasi Pemrosesan Template LaTeX Menggunakan Pendekatan Berbasis AI}
\author{Academic Support Agent}
\date{\today}

\begin{document}
\maketitle

\begin{abstract}

\end{abstract}

\chapter{PENDAHULUAN}


\chapter{TINJAUAN PUSTAKA}


% Original content (commented out):
% \documentclass[12pt]{article}

%
\input{identificacion}

% Carga algunos paquetes y configuraciones particulares
\input{configuracion}

\begin{document}

  \include{portada}

  \tableofcontents

\newpage

  \addcontentsline{toc}{section}{INTRODUCCIÓN}
  \section*{\centering\fontsize{12}{18}\selectfont{\textcolor{blueunadm}{\textbf{INTRODUCCIÓN}}}}
  \input{\rutaintroduccion}

\newpage

  \addcontentsline{toc}{section}{DESARROLLO DE LA ACTIVIDAD}
  \section*{\centering\fontsize{12}{18}\selectfont{\textcolor{blueunadm}{\textbf{DESARROLLO DE LA ACTIVIDAD}}}}
  \input{\rutadesarrollo}

\newpage

  \addcontentsline{toc}{section}{CONCLUSIONES}
  \section*{\centering\fontsize{12}{18}\selectfont{\textcolor{blueunadm}{\textbf{CONCLUSIONES}}}}
  \input{\rutaconclusiones}

\newpage

  \addcontentsline{toc}{section}{REFERENCIAS BIBLIOGRÁFICAS}
  \section*{\centering\fontsize{12}{18}\selectfont{\textcolor{blueunadm}{\textbf{REFERENCIAS BIBLIOGRÁFICAS}}}}
  \printbibliography

\end{document}


\end{document}